\documentclass[12pt]{article}
	\usepackage[utf8]{inputenc}
	\usepackage[T1]{fontenc}
	\usepackage[french]{babel}
	\usepackage{amsmath, amssymb}
	
\begin{document}
\title{L'Équation Normale en Régression Linéaire}
\author{ME ZENAGUI}
\date{\today}
\maketitle

\begin{abstract}
Ce document présente la théorie de l'équation normale appliquée à la régression linéaire. 
Nous introduisons le modèle, dérivons l'équation normale à partir de la minimisation de l'erreur quadratique, 
et discutons ses avantages, limites et variantes.
\end{abstract}

\tableofcontents
\newpage
\section{Introduction}
"Le Machine Learning est un processus qui permet à la machine d'apprendre automatiquement sans être explicitement programmée pour cela." \textit{(Arthur Samuel, 1959)}


Les algorithmes du machine learning se catégorisent principalement en:
\begin{itemize}
    \item \textbf{Apprentissage Supervisé} : L’algorithme dispose d’observations étiquetées.
    \item \textbf{Apprentissage Non Supervisé} : L’algorithme dispose d’observations non étiquetées.
    \item \textbf{Apprentissage Semi-Supervisé} : L’algorithme dispose d’observations partiellement étiquetées.
\end{itemize}


Parmi les principaux problèmes de l’apprentissage supervisé, on distingue : la classification, lorsque l’espace des cibles est dénombrable, et la régression, lorsque l’espace des cibles est continu.

 
La \textbf{régression linéaire} est un modèle statistique utilisé dans l'Apprentissage Machine qui a pour objectif de décrire une variable cible Y par une combinaison linéaire de variables explicatives $X_1, ..., X_n$ \footnote{Source : $https://fortierq.github.io/nb/regression_lineaire$}. 


Dans ce document, nous traiterons l’une des méthodes permettant de résoudre le problème de régression linéaire (voire polynomiale), à savoir adéquation normale. 

\section{Le modèle}


\end{document}